
\chapter{Discussion}
\epigraph{''To me programming is more than an important practical art. It is also a gigantic undertaking in the foundations of knowledge.''}{\vspace{10pt}Grace Hopper}
\label{chapter:discussion}

\newpage

This chapter presents the analysis and discussion of the results shown in chapter \nameref{chapter:results}.

\section{Summary of Findings}

Table \ref{tab:testing_model_global_metrics} shows that across all six configurations spanning three model sizes (Small/Base/Large) and the two heads (with vs. without changes), top‑1 accuracy remained tightly clustered between $0.923$ and $0.928$, indicating that head choice does not materially affect overall accuracy at the dataset level. However, macro‑F1 and per‑class scores reveal important differences. The enhanced head improves macro F1-Score for the Small model ($+0.0048$) while slightly reducing it for the Base ($-0.0016$) and Large ($-0.0035$) models. 

At the class level, as shown in tables \ref{tab:class_base_model_comparison} and 
\ref{tab:class_small_model_comparison} the enhanced head provides consistent gains for Fear emotion class on the F1-Score metric for the Base ($+0.0368$) and Small ($+0.0331$) models's sizes, something similar occurs for Happy emotion class where there is a visible improvement across all sizes. 

The head changes also introduces small regressions for certain majority or easier classes (e.g., Neutral, Angry in Base/Large model sises). These results suggest the enhanced head is most beneficial when backbone capacity is limited (Small/Base)  and less so when the fused backbone representation is already rich, which is the case for the Large model, as reported in table \ref{tab:class_large_model_comparison}.


\section{Error patterns and class imbalance}


%Improvements concentrate on classes that are traditionally considered difficult in FER—such as Fear, which are also among the least represented categories in RAF-DB (Table \ref{tab:class-distribution}). The enhanced head lifts Fear F1-Score in Small and Base models ($+0.0331$ and $+0.0368$), these results can be validated on tables \ref{tab:class_base_model_comparison} and 
%\ref{tab:class_small_model_comparison}. In contrast, classes with abundant, distinctive samples (e.g., Happy) tend to maintain their performance or exhibit small positive shifts. This pattern suggests that the additional non‑linearity and normalization introduced in the enhanced head help the model learn finer decision boundaries low‑prevalence expressions. Meanwhile, overall accuracy remains effectively unchanged because it is dominated by majority classes: as shown in table \ref{tab:class-distribution}, Happy and Surprise represent $38.62 %  $ and $22.16 %$ of the test set, respectively, so small improvements in minority/ambiguous classes barely move the aggregate metrics.
% 	

%Therefore, macro averaged metrics are better aligned with the objective of lifting underperforming categories, and they should be emphasized in the main results table. 

\subsection{Model's Size}

